\section{Why the harmonic value is independent of the formal chart}

% BEGIN G76J1
\paragraph{G76J1. The ambiguity of a normalized exponential frame.}
Work at one framed point over any commutative algebra $S$ of the
specified coefficient field $K$. Give the formal group its logarithm
and use the prescribed Lie eigenbasis for the additive coordinates
$z_i$ (CNC2, CNC7). Its exponential chart is unique. Indeed a change
between two additive coordinates is an additive formal series, hence
linear in characteristic zero, and its tangent map is the identity.
The comparison to these coordinates takes the germ of a section in a
truncation-compatible frame to its ordinary exponential-coordinate germ
with the prescribed gauge correction. This follows by pulling the frame
and its section through the substitution on each truncated polynomial
ring. The coefficient of a monomial of degree $m$ uses only the
substitution and section modulo terms of degree $m+1$.

For two normalized exponential frames their ratios in the truncated
rings are compatible units. Their inverse limit $R$ has constant term
one by rigidification. Write $[b]$ for the rescaling of the
$i$-th coordinate by the embedded eigenvalue of $b$. Comparing the
two normalized line isomorphisms in the adjoint cocycle gives
\[
 [1+b\bar b]R=([b]R)^2 R\,[b\bar b]R .
\]
Indeed the rigidified line isomorphism in this cocycle is unique
(CNC1), so its scalar cancels on taking the ratio. Set $q=\log R$.
In degree $d$, setting $b=1$ gives $(2^d-4)q_d=0$. Thus $q$ is
quadratic. If $c z_i z_j$ is a term, let $s,t$ denote the two CM
embeddings. Taking its coefficient in the logarithmic equation gives
\[
 c\{s(b)s(\bar b)+t(b)t(\bar b)-2s(b)t(b)\}=0.
\]
If $s,t$ differ on the maximal real subfield, an integral real
$b$ separating them gives $c(s(b)-t(b))^2=0$. The nonzero
difference comes from $K$ and is a unit in $S$. Otherwise $t$ is
$s$ or $\bar s$. In the former case choose a nonzero integral
$b$ with $\bar b=-b$ (clear denominators in $a-\bar a$);
then $-4s(b)^2c=0$. Only conjugate pairs can survive. These
arguments also apply over a nonreduced $S$. A section of the $k$-th
line power has its scalar germ multiplied by $R^{-k}=\exp(-kq)$
when the generator is multiplied by $R$.
% END G76J1

% BEGIN G76J2
\paragraph{G76J2. Why the quadratic ambiguity disappears.}
Write the $n$ column germs in disjoint variables $z_{bi}$ and let
$P$ be the original harmonic polynomial in the dual variables
$u_{bi}$. In exponential coordinates the evaluation pairs their product
$G$ with $P$ by $\langle P,G\rangle=P(\partial)G(0)$. For monomials
of multi-index $m$ this is $m!$ times the matching coefficient. In
particular multiplication on the germ side by $z_{bi}$ is adjoint to
$\partial/\partial u_{bi}$ on the polynomial side. Put
\[
 Q=\sum_{b=1}^n q(z_b).
\]
For a monomial $z_i z_j$ of $q$, the adjoint of multiplication by
$\sum_b z_{bi}z_{bj}$ is the contraction
$\sum_b\partial_{u_{bi}}\partial_{u_{bj}}$. By G76J1 only
conjugate embedding pairs enter, and each such contraction annihilates
$P$ by the defining harmonic equations. Thus
$\langle P,Q^hG\rangle=0$ for every $h>0$. Expanding the exponential
gives $\langle P,e^{-kQ}G\rangle=\langle P,G\rangle$.
Only degrees up to $\deg P$ are needed, so the calculation works
unchanged when the equality of germs is given only through that order.
% END G76J2

% BEGIN G76J3
\paragraph{G76J3. Passing from arbitrary frames to compatible ones.}
An original chart specifies a frame at each jet order. Its definition
allows the restriction of one such frame to differ from the preceding
one. Here is how the finite-order argument covers those charts as well.
On the two formal group arguments, the ratio between consecutive frames
restricts to one on both axes. Dividing their laws for addition,
interchange and the adjoint CM action makes its logarithm additive
in each logarithmic coordinate through the given order, symmetric,
and balanced for the adjoint action. Over a rational algebra a
biadditive series is bilinear (differentiate the two addition laws).
Equivalently the ratio is the truncated exponential of a symmetric
bilinear form $B_m$ in the two logarithms. At orders below two this
ratio is one. Choosing compatible bilinear corrections recursively,
with correction zero at order zero and subtracting the successive
defect at each step, makes the corrected frames restrict exactly.
The corrections preserve the axes, symmetry, rigidification and
adjoint balance by the same bilinear identities.

On eigenvectors of types $s,t$, balance says
$(s(a)-t(\bar a))B_m(v_s,v_t)=0$ for every $a\in\mathcal O_E$.
Different embeddings can be separated by an integral $a$, whose
nonzero embedded difference is a unit. The correction therefore pairs
only conjugate types. Restricting to the line frame uses the two
arguments $z,[w_0]z$, and changes the section germ at order $m$
by the exponential of a quadratic form with that same support
(the second eigenvector acquires its eigenvalue of $w_0$).
For each of the given sections, express its pullback in each corrected
jet frame. Uniqueness of the coefficient of a generator makes these
coefficients compatible under restriction and hence gives the corrected
germ. Only the germs specified in the original comparison are
required in the two starting charts. Use $m=\deg P$ in G76J2. The correction changes no harmonic value;
there is no need to identify the separate corrections $B_m$ as one
matrix. Compare the two corrected charts in their common exponential
coordinates by G76J1--G76J2. This proves equality of the original
harmonic evaluations at every framed point, for every such pair of
charts and every tuple of germs of the same sections. In particular
the harmonic expression built from any two connections satisfying
the prescribed chart formula has the same value, by inserting
their respective germs from CNC15.
% END G76J3
